\documentclass[11pt]{article}

\usepackage[margin=1.15in]{geometry}
\usepackage{booktabs}
\usepackage{tabularx}
\usepackage{xcolor}
\usepackage{tikz}
\usepackage{amsmath}
\usepackage{amssymb}
\usepackage{amsthm}
\usepackage{microtype}
\usepackage{natbib}
\usepackage{enumitem}
\usepackage[colorlinks=true,linkcolor=paperblue,citecolor=paperblue,urlcolor=paperblue]{hyperref}

\definecolor{paperblue}{RGB}{42,79,168}
\definecolor{rulegray}{RGB}{140,146,158}
\definecolor{washblue}{RGB}{232,237,248}

\usetikzlibrary{positioning,arrows.meta,fit,backgrounds}

\bibliographystyle{plainnat}
\setlength{\bibsep}{4pt}
\setlist{itemsep=2pt,topsep=4pt}

\theoremstyle{definition}
\newtheorem{definition}{Definition}

\newcommand{\sut}{\ensuremath{S}}
\newcommand{\term}[1]{\emph{#1}}

\title{\vspace{-2.2em}\textbf{Release Gates for Physical AI}\\[0.35em]
\large A Unified Framework for Functional, Adversarial, Cyber, and Hardware\\
Verification of Learned Robot Systems}

\author{Shubh Jain\\ \small Carnegie Mellon University \\ \small\texttt{shubhj@andrew.cmu.edu}}

\date{\small Preprint v0.1.0 \quad\textbullet\quad 30 September 2026 \quad\textbullet\quad Not peer reviewed}

\begin{document}
\maketitle

\begin{center}
\fcolorbox{rulegray}{washblue}{%
\begin{minipage}{0.92\textwidth}
\small
\textbf{Status.} This is a framework and position paper. It contains no new experimental
results. Every quantitative figure is attributed to published work or to a named standard;
every threshold proposed here is labelled as a proposal and is not yet validated.
\textbf{Scope.} Manipulation and station-based mobile manipulation. Locomotion, whole-body
balance, and public-road autonomy are out of scope.
\end{minipage}}
\end{center}

\begin{abstract}
\noindent Learned manipulation policies are now updated on the cadence of software, shipped
to fleets monthly, and evaluated with the rigour of a demo. Classical machine-safety
verification assumes a deterministic controller whose behaviour can be argued from its
structure; a vision-language-action policy offers neither determinism nor structure, and its
competence is a property of a distribution rather than of a program. The decision to deploy
a new policy version is consequently made today without a coherent body of evidence.

We observe that the evidence bearing on that decision already exists in four separate
literatures and four separate industrial practices---functional evaluation, adversarial
robustness, product cybersecurity, and hardware reliability---and that no framework
composes them. We propose one. Our contributions are: (i) a taxonomy of the four test
families that bear on a release decision, with the metrics each can produce; (ii) a
\term{fidelity-gating} principle, under which a test may block a release only once its
predictive validity against real outcomes has been measured, and otherwise reports as
advisory; (iii) a release-scoped, immutable evidence model mapped to the clauses of
ISO~10218:2025 and the EU Cyber Resilience Act, including the treatment of model weights as
components of a software bill of materials; and (iv) an evaluation protocol for the
framework itself, centred on a retrospective study that asks whether a proposed gate would
have caught regressions that actually reached the field.

We further argue that the composition is not merely administrative: several failure modes
are visible only across families---a policy version that improves task success while
becoming more susceptible to in-scene prompt injection, or whose grasp failures concentrate
on units whose end-effectors are past a wear threshold---and are structurally invisible to
any single-family evaluation.
\end{abstract}

\section{Introduction}

A robot deployment running a learned policy carries two independent kinds of uncertainty:
whether the policy will do the task, and whether the surrounding system will behave when
something adversarial, degraded, or unforeseen occurs. Industrial practice has mature
answers for neither, and no answer at all for their interaction.

The immediate cause is cadence. A fleet running a vision-language-action (VLA) policy
receives new weights monthly or faster. The governing safety standard for industrial robot
applications was revised in 2025 to accommodate, among other things, machinery incorporating
machine-learning functions that alter behaviour over time
\citep{iso102181,iso102182}, but conformity assessment remains an event rather than a
subscription. Between two certified states a policy may change behaviour substantially, and
nothing in the classical process observes that change.

The deeper cause is epistemic. Functional safety argues from structure: a speed-and-separation
monitor can be shown to be correct. A learned policy has no structure to argue from, so its
competence must be \term{measured}, on a distribution, repeatedly. Measurement of that kind
produces evidence with an unfamiliar property: it is only as good as its own predictive
validity. A simulated benchmark that does not rank policy versions the way reality ranks
them is not weak evidence but misleading evidence, and a release process that treats the two
alike is worse than one with no simulation at all.

\subsection{The four silos}

Ask a robot company for the evidence behind its last release and four disconnected artifacts
arrive: a notebook of task success rates on held-out cells; a penetration-test report,
usually more than a year old, that does not mention the policy; maintenance records in a
separate system describing hardware whose degradation nobody has related to task
performance; and a safety file at a notified body describing a configuration that predates
the current weights.

Each silo has a serious literature. Real-to-sim evaluation now reports rank correlations
against real-robot outcomes of $0.911$ \citep{simfoundry} and $0.887$ \citep{robosnap}.
Physical prompt injection has a systematic taxonomy and a 5{,}670-trial benchmark
\citep{paperhijack}. Product cybersecurity has a regulation with dated obligations and
defined artifacts \citep{cra}. Hardware reliability has an established engineering practice.
What is missing is the composition, and the composition is where the release decision lives.

\subsection{Contributions}

\begin{enumerate}[leftmargin=1.4em]
  \item A \textbf{test taxonomy for the release decision} (\S\ref{sec:families}): four
  families, their runners, and the metrics each can produce with current published methods.
  \item \textbf{Fidelity gating} (\S\ref{sec:gating}): a rule separating tests permitted to
  block a release from tests permitted only to advise, based on measured predictive validity
  rather than on apparent realism.
  \item A \textbf{release-scoped evidence model} (\S\ref{sec:store}--\S\ref{sec:compliance})
  with an explicit mapping to ISO~10218:2025 and CRA artifacts, including model weights as
  SBOM components.
  \item An \textbf{evaluation protocol for the framework} (\S\ref{sec:protocol}), centred on
  a retrospective study of known field regressions.
  \item A statement of \textbf{open problems} (\S\ref{sec:open}), the most consequential
  being that the verifiers this framework depends on are, on contact-rich tasks, close to
  unusable today \citep{failbench}.
\end{enumerate}

\section{Related work}

\subsection{Functional evaluation}

Evaluating a policy in simulation against real outcomes has moved from aspiration to
measurement. SimplerEnv established visual-matching and variant-aggregation protocols for
comparing simulated and real manipulation performance \citep{simplerenv}. Recent automated
real-to-sim systems report correlation directly: SimFoundry builds simulation-ready scenes
from RGB video and reports a mean Pearson correlation of $0.911$ with a maximum rank
violation of $0.018$ across seven tasks and five policy types, at roughly five minutes of
reconstruction per object, with $F_1$ between $0.81$ and $0.92$ fully automated and
$0.93$--$0.99$ given about three minutes of per-object human refinement \citep{simfoundry}.
RoboSnap reconstructs from a single image and reports a Pearson correlation of $0.887$ and
successful replay of five of five real trajectories \citep{robosnap}. Both are explicit that
physical parameters---mass, friction---are assigned by querying a vision-language model
rather than measured; RoboSnap states that it includes no dedicated pipeline for automatic
physical-parameter estimation \citep{robosnap}, and an earlier framework attributes its
residual error to ``domain mismatch, mesh reconstruction errors, and discrepancies in
material property estimates'' \citep{scanmatsim}.

A second line replaces the simulator with a learned world model and scores policies inside
it \citep{onexwm,gigaworld,dworldeval}. The trade is legibility: a learned evaluator yields
no contact forces, no queryable state, and no counterfactual outside its training
distribution, which makes it unsuitable as the substrate for cyber or hardware evidence even
where its functional predictions are good.

Underneath both sits an unresolved dependency: deciding whether an episode succeeded.
FailBench evaluates thirteen VLM-based success detectors and finds the best reaching $0.77$
mean balanced accuracy, degrading to below $0.60$---near chance---on contact-intensive
assembly, with a systematic bias toward predicting success under ambiguous evidence; models
fine-tuned for failure detection underperformed general-purpose VLMs \citep{failbench}. A
companion survey catalogues verifier designs and their failure modes \citep{nofreechecker},
and calibrated task-success confidence is an active target \citep{vlaconf}.

\subsection{Adversarial robustness}

Embodiment changes the adversarial threat model, because the consequence is physical and
frequently irreversible \citep{vlasafety}. The sharpest current result is physical prompt
injection: adversarial text placed in the robot's visual field acts as an indirect injection
into a VLM planner's reasoning. Across 5{,}670 trials with a four-category
taxonomy---indirect signage, task redefinition, authority impersonation, conflict
injection---attack success reached $27.0\%$, $29.4\%$, and $5.0\%$ on three frontier models
\citep{paperhijack}. Related work demonstrates trajectory-level redirection, in which the
instruction still appears to specify the intended task while the physical outcome is changed
\citep{trajredirect}, partially observable adversarial patches \citep{patchattack}, and
interference through physical signals such as sound, light, and electromagnetic emission.

Defensively, adversarial scenario generation is being formulated as a red-team/blue-team
game \citep{advsynth}, boundary-focused sampling is used for failure discovery before
deployment \citep{robogate}, and the formal-methods community has surveyed falsification for
learned policies \citep{formalsurvey}. None of this work is packaged as release evidence.

\subsection{Cybersecurity and standards}

Two instruments set dates. The revised ISO~10218-1 and 10218-2, published in 2025, absorb
collaborative requirements formerly in ISO/TS~15066 and add robot classifications with
matching functional-safety requirements, safety-related cybersecurity requirements, and
end-effector guidance \citep{iso102181,iso102182}. ISO~25785-1, the first standard written
for dynamically stable industrial mobile robots---the humanoids and quadrupeds now entering
warehouses---remained a committee draft in early 2026 \citep{iso25785}. Separately, the EU
Cyber Resilience Act entered into force in December 2024; vulnerability and incident
reporting obligations apply from 11 September 2026, and full compliance---CE marking,
machine-readable SBOM, secure-by-design requirements, vulnerability handling, and security
updates across the support period---from 11 December 2027, with penalties up to
EUR~15M or $2.5\%$ of global turnover \citep{cra,cratimeline}.

Vendor safety stacks have moved faster than the standards. Halos for Robotics, announced in
June 2026, couples safety compute, sensor bridging, a safety OS, and an accredited
inspection lab that prepares integrations for third-party certification against IEC~61508,
ISO~13849, and ISO/IEC~TR~5469 \citep{halos}; FORT Robotics extends it with
infrastructure-side sensing \citep{fort}. This is functional safety: it bounds harm to
people, does not measure task competence, and does not address adversarial manipulation of
the policy's inputs. Control-barrier formulations that embed ISO~10218 constraints directly
in the controller \citep{cbfiso} share that boundary.

\subsection{Hardware reliability and its coupling to policy}

Industrial practitioners consistently identify actuation as the binding constraint: a
full-body platform carries thirty or more degrees of freedom, each required to deliver
repeatable torque across thousands of hours, and harmonic-drive flex-splines are subject to
fatigue under exactly the cyclic loading that walking and repetitive assembly produce;
end-effectors are reported as the most failure-prone and most expensive subsystem
\citep{humanoidhw,humanoidchal}. Absent from the literature is the coupling: no published
framework relates a measured degradation state---a suction cup past a wear threshold, a
drifting camera extrinsic---to the task success of the specific policy version running on
that unit. In deployment this coupling is routine, and it is the most common informal
explanation for ``the same policy performs differently across sites.''

\subsection{What no existing system composes}

Independent evaluation infrastructure is emerging---community and hosted real-robot
benchmarks \citep{roboarena,robochallenge}, low-cost reproducible rigs \citep{phail,replica}---and
is converging on public leaderboards for functional performance. Industrial V\&V platforms in
the adjacent automotive domain integrate reconstruction, scenario generation, and
regulatory-grade reporting, with a published three-tier metric framework separating
reconstruction fidelity, open-loop stack agreement, and closed-loop behavioural comparison
against real drive data \citep{neuralsim,wfm}. We take that three-tier structure as prior art
and generalise it, noting that its robotics analogue must additionally carry cyber and
hardware evidence, because in manipulation the same release decision is gated on all four.

\section{Problem statement}
\label{sec:problem}

\begin{definition}[System under test]
A system under test is the tuple
\begin{equation}
  \sut = (h,\; \pi,\; \sigma,\; c)
\end{equation}
with $h$ the hardware configuration (kinematics, end-effector, sensors, and their measured
degradation state), $\pi$ the policy checkpoint, $\sigma$ the software and firmware stack
including model dependencies, and $c$ the site configuration (fixtures, item catalogue,
lighting, network posture). A release changes at least one component; evidence is valid only
for the tuple it was produced against.
\end{definition}

\begin{definition}[Evidence]
An evidence record is
\begin{equation}
  e = (t,\; r,\; \sut,\; y,\; \rho,\; p)
\end{equation}
for test specification $t$, runner $r$, system under test $\sut$, outcome $y$, fidelity
estimate $\rho$ for the pair $(t,r)$, and provenance $p$ (inputs, versions, seeds,
operator). Records are immutable and scoped to a release.
\end{definition}

\begin{definition}[Gate]
For a candidate $\sut'$ and baseline $\sut$, a family of tests partitions into a gating set
$G$ and an advisory set $A$ by
\begin{equation}
  t \in G \iff \rho(t,r) \ge \tau \;\wedge\; n(t) \ge n_{\min}.
\end{equation}
The release is blocked if any $t \in G$ shows a regression beyond its tolerance. Tests in
$A$ are reported with their fidelity and cannot block. We propose $\tau = 0.8$ as a
rank-correlation threshold (\S\ref{sec:metrics}), by analogy with the $0.887$--$0.911$
correlations reported by current real-to-sim systems \citep{simfoundry,robosnap}; the value
requires empirical justification per domain.
\end{definition}

Three consequences follow. First, \textbf{fidelity is a property of a test--runner pair, not
of a simulator}: the same engine may be gating-grade for a rigid pick and advisory-grade for
cable routing. Second, \textbf{evidence expires} when any component of $\sut$ changes, which
is why hardware degradation belongs inside the model rather than beside it. Third, \textbf{a
release decision is a multiple-comparison problem} across tasks and families, and must be
treated as such rather than as a single averaged score.

\section{Framework}

\begin{figure}[t]
\centering
\begin{tikzpicture}[
  font=\sffamily\scriptsize,
  box/.style={draw=rulegray, fill=white, minimum height=7mm, text width=26mm, align=center, inner sep=2pt},
  sut/.style={draw=paperblue, fill=washblue, minimum height=16mm, text width=26mm, align=center},
  store/.style={draw=rulegray, line width=0.8pt, fill=white, minimum height=24mm, text width=28mm, align=center},
  outbox/.style={draw=rulegray, fill=white, minimum height=10mm, text width=27mm, align=center},
  arw/.style={-{Latex[length=1.6mm]}, draw=rulegray}
]
\node[sut] (s) {\textbf{SYSTEM UNDER TEST}\\[2pt] $h \cdot \pi \cdot \sigma \cdot c$};
\node[box, right=9mm of s, yshift=13.5mm] (f1) {Functional};
\node[box, right=9mm of s, yshift=4.5mm]  (f2) {Adversarial};
\node[box, right=9mm of s, yshift=-4.5mm] (f3) {Cyber};
\node[box, right=9mm of s, yshift=-13.5mm](f4) {Hardware};
\node[draw=rulegray, dashed, right=8mm of f2, yshift=-4.5mm, minimum height=30mm, text width=25mm, align=left, inner sep=4pt] (run)
  {calibrated sim\\ log replay\\ hardware-in-loop\\ physical cell\\ static analysis\\[3pt] \textit{each emits} $\rho$};
\node[store, right=8mm of run] (ev) {\textbf{EVIDENCE STORE}\\[3pt] immutable\\ release-scoped\\ fidelity-tagged};
\node[outbox, right=8mm of ev, yshift=8mm] (gate) {\textbf{RELEASE GATE}\\ pass $\cdot$ block $\cdot$ advise};
\node[outbox, right=8mm of ev, yshift=-9mm] (pack) {\textbf{COMPLIANCE PACK}\\ clause-mapped};
\draw[arw] (s.east) -- (f2.west);
\draw[arw] (s.east) -- (f3.west);
\draw[arw] (f1.east) -- ++(4mm,0);
\draw[arw] (f4.east) -- ++(4mm,0);
\draw[arw] (run.east) -- (ev.west);
\draw[arw] (ev.east) -- (gate.west);
\draw[arw] (ev.east) -- (pack.west);
\draw[arw, dashed, draw=paperblue] (pack.south) -- ++(0,-6mm) -| (s.south);
\end{tikzpicture}
\caption{The release-gate pipeline. Four test families execute on five runner types; every
result carries a fidelity estimate that determines whether it may block a release. The
evidence store is the system of record from which both the gate decision and the
clause-mapped compliance pack are derived. Field failures and degradation telemetry re-enter
as new test specifications, scoped to the tuple on which they were observed.}
\label{fig:pipeline}
\end{figure}

\subsection{Test families and their metrics}
\label{sec:families}

Table~\ref{tab:families} states what each family tests, where it can run today, and what it
produces. The \term{status} column records our assessment of whether current published
methods support gating-grade evidence, and is the honest content of this paper: two of four
families cannot presently gate on contact-rich work.

\begin{table}[t]
\small
\centering
\begin{tabularx}{\textwidth}{@{}l X l X@{}}
\toprule
\textbf{Family} & \textbf{Tests} & \textbf{Primary metric} & \textbf{Status} \\
\midrule
Functional & Task success per family, cycle time, cross-task regression, field-failure replay
  & Success-rate delta vs.\ baseline, per task family, with interval
  & Gating for rigid tasks; advisory for deformable and insertion \\
Adversarial & In-scene prompt injection, patches, OOD scenes, sensor interference, trajectory redirection
  & Attack success rate per category, pre/post mitigation
  & Gating: attack success is directly observed, not inferred \\
Cyber & SBOM completeness incl.\ model artifacts, known-vulnerability exposure, update integrity, access paths
  & Unresolved exposure count by severity; time-to-evidence
  & Gating: deterministic, and externally mandated \\
Hardware & Actuator and transmission wear, end-effector degradation, calibration drift, thermal envelope
  & Degradation index and projected service interval
  & Advisory: coupling to task success unmodelled (\S2.4) \\
\bottomrule
\end{tabularx}
\caption{The four families. Status reflects the authors' assessment against the surveyed
literature, not a measured result.}
\label{tab:families}
\end{table}

\subsection{Runners}

A runner is any mechanism that can execute a test specification and return an outcome: a
calibrated simulator, replay of logged episodes against a policy, hardware-in-the-loop, a
physical cell, or static analysis of the software stack. Runners are interchangeable within a
family and are the unit at which fidelity is estimated. The framework deliberately does not
privilege simulation. A test whose only trustworthy runner is a physical cell remains a
first-class test; it is simply expensive, and its sampling budget must be allocated
accordingly (\S\ref{sec:budget}).

\subsection{Fidelity gating}
\label{sec:gating}

A test earns the right to block a release; it does not hold that right by default, and
realism is not the qualification---predictive validity is. Operationally: for each
test--runner pair, estimate the rank correlation between its outcomes and real outcomes
across a set of historical checkpoints of the same system; report it with every result;
permit gating only above $\tau$. This rule matters more in robotics than in software testing
because a robotics test can look extremely convincing and still be uninformative---a
photorealistic scene with guessed friction coefficients produces confident, wrong rankings,
which is precisely the configuration current automated real-to-sim pipelines ship in
\citep{simfoundry,robosnap,scanmatsim}.

Fidelity gating has a second effect that we consider its main practical justification: it
makes low-fidelity tests \term{safe to keep}. A cable-routing suite with $\rho = 0.5$ still
surfaces a large regression worth a real-world check; under a naive gate it would either
block releases spuriously or be deleted. Advisory status retains the signal while denying it
authority.

\subsection{Evidence store}
\label{sec:store}

The store is append-only, addressed by release, and holds provenance sufficient to re-run
each record. Three properties are load-bearing. \textbf{Immutability}, because an audit that
can be edited is not an audit. \textbf{Release scoping}, so a query returns ``what did we
know when we shipped 4.2'' rather than a current-state view. \textbf{Fidelity tagging}, so
the gate rule is computable from the store rather than applied out of band. A machine-readable
schema for the record of Definition~2 accompanies this paper.

\subsection{Compliance mapping}
\label{sec:compliance}

Evidence produced for engineering reasons should discharge regulatory obligations without a
second campaign. Table~\ref{tab:compliance} gives the mapping we propose; the novel element
is the treatment of model weights as SBOM components, which the CRA's machine-readable SBOM
requirement \citep{cra} does not contemplate and which no current tooling emits.

\begin{table}[t]
\small
\centering
\begin{tabularx}{\textwidth}{@{}X l X@{}}
\toprule
\textbf{Obligation} & \textbf{Instrument} & \textbf{Evidence that discharges it} \\
\midrule
Vulnerability and incident reporting & CRA, from 11 Sep 2026 & Cyber-family exposure records with timestamps and disposition \\
Machine-readable SBOM & CRA, full 11 Dec 2027 & Stack manifest extended with model artifacts: weights hash, training-data provenance class, licence \\
Security updates over support period & CRA & Release history with per-release gate outcome and update integrity test \\
ML functions altering behaviour over time & ISO 10218-1:2025 & Functional-family regression record per release, plus adversarial-family results \\
Application and cell requirements & ISO 10218-2:2025 & Site-scoped evidence bound to $c$ in Definition~1 \\
Dynamically stable mobile robots & ISO 25785-1 (draft) & Out of scope pending publication; structure anticipated \\
\bottomrule
\end{tabularx}
\caption{Proposed mapping from engineering evidence to regulatory obligation. We produce
evidence; we do not certify. Certification remains with accredited bodies, and vendor safety
stacks occupy the functional-safety layer beneath this table.}
\label{tab:compliance}
\end{table}

\subsection{Cross-family inference}
\label{sec:cross}

The composition claim is empirical and testable. We predict three interaction effects that no
single-family evaluation can express.

\begin{enumerate}[leftmargin=1.4em]
  \item \textbf{Robustness--competence trade-off.} A checkpoint that improves functional
  success may increase susceptibility to in-scene injection, if the gain came from stronger
  instruction-following. Given the per-model spread reported in \citep{paperhijack}---$27.0\%$,
  $29.4\%$, $5.0\%$---variation of this magnitude across checkpoints of a single model family
  would be unsurprising and is, to our knowledge, unmeasured.
  \item \textbf{Degradation-conditioned failure.} Functional failures should concentrate on
  units whose hardware degradation index exceeds a threshold. Stratifying functional results
  by hardware state converts unexplained cross-site variance into a maintenance action.
  \item \textbf{Fidelity--family asymmetry.} A runner's fidelity is family-dependent: a
  simulator adequate for adversarial scene perturbation may be inadequate for the contact
  physics of the same task. The framework records fidelity per pair precisely so this
  asymmetry is representable.
\end{enumerate}

\section{Metrics and proposed thresholds}
\label{sec:metrics}

Table~\ref{tab:metrics} collects the metrics the framework requires. Thresholds marked
\term{proposed} are ours and unvalidated; they are stated numerically so they can be
falsified rather than debated.

\begin{table}[t]
\small
\centering
\begin{tabularx}{\textwidth}{@{}X r X@{}}
\toprule
\textbf{Metric} & \textbf{Threshold} & \textbf{Basis} \\
\midrule
Suite fidelity $\rho$: rank correlation between suite and real outcomes over $\ge 8$ historical checkpoints & $\ge 0.80$ & Proposed; anchored to $0.887$--$0.911$ \citep{simfoundry,robosnap} \\
Rank violation: maximum inversion between suite and real ranking & $\le 0.05$ & Proposed; $0.018$ in \citep{simfoundry}, $0.0066$ in \citep{robosnap} \\
Verifier balanced accuracy, reported separately for contact-light and contact-heavy tasks & $\ge 0.90$ & Proposed; best published $0.77$ overall, $<0.60$ contact-intensive \citep{failbench} \\
Attack success rate per category, pre/post mitigation & report & Measured directly; protocol from \citep{paperhijack} \\
Regression catch rate on known past field regressions & $\ge 0.80$ & Proposed; \S\ref{sec:retro} would establish it \\
Time-to-evidence: checkpoint availability to complete gate decision & $\le 24$\,h & Proposed; set by release cadence \\
Test-time compute cost as multiple of base inference & report & Best-of-$N$ at $\approx$2.2\,Hz \citep{cover}; adaptive matches best-of-10 at $34\%$ lower latency \citep{elastic} \\
\bottomrule
\end{tabularx}
\caption{Metrics, with thresholds stated as falsifiable proposals rather than findings.}
\label{tab:metrics}
\end{table}

\section{Evaluation protocol for the framework}
\label{sec:protocol}

A framework of this kind is validated by whether its gate decisions would have been right. We
specify four studies in increasing cost.

\subsection{Retrospective gate study}
\label{sec:retro}
Take a deployed fleet's last $N$ policy checkpoints, their real per-task outcomes, and their
known field regressions. Construct suites without access to the outcomes, compute $\rho$ per
suite, apply the gate rule, and report rank correlation, regression catch rate, and
false-block rate. This is the load-bearing experiment; it requires an industrial partner and
no new capability.

\subsection{Adversarial pre/post study}
Instantiate the four attack categories of \citep{paperhijack} against a fixed policy, measure
attack success, apply a guard layer, and re-measure. Report transfer from simulated to
physical presentation of the same attack, which is the quantity deciding whether simulated
adversarial evidence may gate.

\subsection{Cross-family interaction study}
For a sequence of checkpoints, measure functional success and attack success jointly, and
stratify functional failures by hardware degradation index. The predictions of
\S\ref{sec:cross} are either present or absent; either result is informative, and neither has
been reported.

\subsection{Budget allocation study}
\label{sec:budget}
Given a fixed budget of physical-cell hours, compare allocation policies---uniform,
fidelity-weighted, uncertainty-weighted---on regression catch rate. The relevant baseline is
random sampling, which falsification work suggests is weak but which industrial practice
currently approximates \citep{robogate,formalsurvey}.

\section{Open problems}
\label{sec:open}

\begin{enumerate}[leftmargin=1.4em]
  \item \textbf{Verifiers on contact-rich tasks.} The functional leg rests on success
  detection, which is near chance exactly where industrial value concentrates
  \citep{failbench}. We believe the missing evidence is non-visual---wrist force/torque,
  joint effort, vacuum pressure, tactile---and that verifier design should be multimodal
  rather than vision-only. This is a claim, not a result.
  \item \textbf{Calibrating fidelity cheaply.} Estimating $\rho$ requires historical
  checkpoints with real outcomes; new deployments have none. Transfer of fidelity estimates
  across sites and embodiments is unsolved.
  \item \textbf{Physical parameters.} Every automated real-to-sim pipeline surveyed obtains
  mass and friction from a language model's prior \citep{simfoundry,robosnap,scanmatsim}.
  Identification from interaction logs is the obvious alternative and only partially explored
  \citep{twinaligner,asid}.
  \item \textbf{Compute budget of verification.} Verification-time scaling improves outcomes
  substantially but currently at control rates near 2.2\,Hz \citep{cover}; adaptive
  allocation is promising \citep{elastic} but trains per task.
  \item \textbf{Standards coverage for learned components.} ISO~10218:2025 acknowledges ML
  functions but does not specify acceptance evidence for them; ISO~25785-1 is unpublished
  \citep{iso25785}. That gap is where this framework is intended to be useful, and also why
  it cannot yet claim regulatory standing.
  \item \textbf{Privacy as an evidence constraint.} First-person and workplace video is
  treated as biometric data under GDPR, with anonymisation expected before a dataset leaves
  the collection site. Evidence sharing between deployer, model provider, and assessor is
  therefore a data-protection design problem, not only a technical one.
\end{enumerate}

\section{Limitations}

This paper proposes structure and measures nothing. Its thresholds are guesses with stated
provenance. The status assignments in Table~\ref{tab:families} are judgements formed from the
surveyed literature and would shift with a single strong result on contact-rich verification.
Several industrial figures we cite---actuator fatigue as the binding constraint,
end-effectors as the most failure-prone subsystem---come from trade and vendor sources rather
than peer-reviewed work and should be read as practitioner consensus rather than measurement.
The framework also assumes a release is a discrete event; continuous or per-site online
adaptation breaks that assumption, and we do not address it.

\section{Conclusion}

Robotics does not lack tests. It lacks a rule for which tests are allowed to stop a release,
and a place to keep the answer. The four families already exist in the literature; what the
field has not done is compose them, tag each result with its measured predictive validity,
and hold the composition against the release decision it is meant to inform. We have
specified that composition precisely enough to be implemented and argued against, and named
the one dependency that would make it work or fail: a verifier that can tell, on
contact-rich tasks, whether the robot actually succeeded.

\subsection*{Reproducibility and artifacts}
The project page, the LaTeX source of this paper, and a JSON Schema for the evidence record of
Definition~2 are available in the accompanying repository. No datasets or trained models are
released, as none were produced.

\bibliography{refs}

\end{document}
